\documentclass[11pt]{article}
\usepackage{amsmath,amssymb,mathtools}
\usepackage{fontspec}
\usepackage{unicode-math}
\setmainfont{Latin Modern Roman}
\setmathfont{Latin Modern Math}
\usepackage{xurl}
\usepackage[unicode,colorlinks=true,linkcolor=blue,urlcolor=blue,pdftitle={Problem statement}]{hyperref}
\newcommand{\sref}[2]{\hyperlink{source:#1:#2}{[S#2]}}
\newcommand{\cataloguescope}[1]{\paragraph{Mathematical statement.}#1}
\begin{document}
\section{Research Problems in Function Theory — Problem 2.18}\label{statement}

\subsection{Definitions and mathematical statement}

\paragraph{Shared notation.}$\mathbb N=\{1,2,\ldots\}$, and $\mathbb Z,\mathbb Q,\mathbb R,\mathbb C$ denote the integers, rationals, reals and complex numbers. Any other conventions are specified locally.


For an entire function $f$, write
\[
 M(r,f)=\max_{|z|=r}|f(z)|,\qquad m_0(r,f)=\min_{|z|=r}|f(z)|.
\]
For nonconstant $f$, its order is $\rho(f)=\limsup_{r\to\infty}\log\log M(r,f)/\log r$;
the lower order is defined using $\liminf$ instead. Every constant entire function, including the zero function, has order and lower order zero. In logarithmic modulus ratios, use the extended-real convention $\log0=-\infty$.
For a function holomorphic on a neighbourhood of a circle $|z|=r_0>0$, define $M(r,f)$ for nearby radii and, where defined,
$b_f(r)=(r\,d/dr)^2\log M(r,f)$. Write $b_f(r_0-0)$ and $b_f(r_0+0)$ for the left and right limiting values.
Exclude monomials as required by Update 2.17, which Update 2.18 also invokes.

\paragraph{Statement recorded in the supplied source.}
Can a nonmonomial entire function satisfy
\[b_f(r_0-0)=b_f(r_0+0)=0\]
at some $r_0>0$? More generally, can this happen for a function holomorphic on a neighbourhood of $|z|=r_0$? \sref{statement}{2}
Update 2.18 gives a positive answer due to London, using an entire function $(1-z)\exp(\alpha z+\beta z^2+\gamma z^3)$. It distinguishes this from vanishing throughout an interval of radii. \sref{statement}{2}
\subsection{Short English statement}
Can both one-sided curvatures of the logarithmic maximum modulus vanish at the same circle, without the function being a monomial?
\subsection{Sources}
\begin{description}
\item[\hypertarget{source:statement:1}{S1}] ULAM.AI and the UnsolvedMath Contributors, \emph{UnsolvedMath: A Curated Collection of Open Mathematics Problems}, catalogue entry. CC BY 4.0. \url{https://huggingface.co/datasets/ulamai/UnsolvedMath}.
\item[\hypertarget{source:statement:2}{S2}] W. K. Hayman and E. F. Lingham, \emph{Research Problems in Function Theory} (2018), Chapter 2, Problem 2.18 and Update 2.18. \url{https://arxiv.org/abs/1809.07200}.
\end{description}
\label{end:statement}
\end{document}
